% =====================================================================
% intro_conc_outline.tex
% Draft Introduction (Chapter 1) and Conclusion (final chapter). They sit
% at opposite ends of the report; kept in one file for drafting. Written in
% the established style (British spelling, \texttt{} library names, energy-
% literal wording, \S\ref cross-refs). "% TODO" marks anything to verify,
% cite, or calibrate.
%
% Energy-framing reminder: literal "lowest/highest energy", "uses the
% least/most energy". No speed/efficiency words (faster, efficient,
% cheapest, performance-as-speed).
% =====================================================================

\chapter{Introduction}
\label{chap:Introduction}

The energy a program draws is increasingly treated as a property to be
measured and managed in its own right, rather than as an incidental
consequence of how long the program runs. This report studies that property
for a narrow but pervasive task in the .NET ecosystem: the serialisation and
deserialisation of JSON.

\section{Motivation}
\label{sec:intro-motivation}

The electricity consumed by data centres has grown alongside the workloads
they host, and the resulting operational cost and carbon footprint have made
software energy consumption a practical engineering concern as much as an
environmental one. % TODO cite: software energy / green computing / data-centre energy
Optimisation effort has traditionally targeted execution time, on the
assumption that reducing how long a program runs also reduces the energy it
draws. Energy, however, depends on how the underlying hardware is exercised,
and treating it as a distinct quantity allows that assumption to be tested
rather than presumed.

Within this broad concern, JSON serialisation is a task worth singling out.
JSON is among the most widely used data-interchange formats, and the
conversion between in-memory objects and their textual form sits on the hot
paths of most networked .NET applications, including web APIs, inter-service
messaging, configuration loading, caching, and persistence. % TODO cite: JSON ubiquity
Because this conversion is performed constantly and at scale, even a small
per-operation difference in energy is repeated across the lifetime of a
long-running service, where it accumulates into a cost worth attributing to
the choice of library.

The .NET ecosystem offers several mature JSON libraries that approach the same
task in materially different ways. \texttt{Newtonsoft.Json} is the
long-established, feature-rich option; \texttt{System.Text.Json}, Microsoft's
modern library, can operate either through runtime reflection
(\texttt{STJRefGen}) or through a compile-time source generator
(\texttt{STJSrcGen}); and \texttt{SpanJson} and \texttt{Utf8Json} are built
around UTF-8 byte buffers, with \texttt{SpanJson} additionally using
SIMD-accelerated parsing. These libraries differ in how they allocate memory,
how they encode and decode text, and how much work they move from run time to
compile time. Each of these choices plausibly affects the energy a library
draws, and not necessarily in proportion to the time it takes.

\section{Problem Statement}
\label{sec:intro-problem}

This report extends the authors' 9th-semester
project~\cite{geotoms2025microbenchmarking}, which integrated \gls{rapl}
energy readings into BenchmarkDotNet and compared five .NET JSON libraries
across serialisation and deserialisation. That study established that the
libraries differ substantially in energy and that their ranking depends on
payload size, but it measured only three real-world GitHub API documents in a
single clean environment. Three documents conflate the properties that
distinguish JSON workloads --- size, nesting depth, structural width, and
content type vary at once --- so they cannot reveal which property drives a
given energy difference, nor whether a ranking observed on them generalises.
They also leave unexamined whether such a ranking survives outside a quiet
laboratory. More broadly, public comparisons of these libraries typically
report throughput or latency, from which energy is inferred rather than
measured directly. % TODO cite: existing JSON benchmarks / throughput comparisons

This leaves two questions open. The first concerns the workload: how does a
library's energy use depend on the shape of the data it processes --- the
number of objects, the nesting depth, the structural width, and the type and
length of the values --- and does the energy ranking between libraries hold
across that space, or does the lowest-energy library change with the
workload? A ranking established on one kind of document need not transfer to
another, and without sweeping the workload systematically there is no way to
separate the robust differences from the artefacts of a single input.

The second question concerns the measurement itself. Energy readings are
sensitive to the machine on which they are taken and to whatever else that
machine is doing, so a ranking obtained in a quiet laboratory may not survive
contact with a loaded, production-like environment. For a measurement
methodology to be useful, it must produce rankings that remain consistent
under such conditions. This report therefore evaluates the
\gls{metrion-profiler}, a profiler under active development, against the
hardware \gls{rapl} interface, and tests whether its rankings remain stable
when the environment is deliberately stressed.

\section{Research Questions}
\label{sec:intro-rqs}

These concerns motivate the three research questions that this report
addresses.

\begin{description}
  \item[RQ1] How do libraries' energy rankings and scaling rates differ
    across workload dimensions?
  \item[RQ2] What architectural and runtime factors explain the observed
    differences between libraries?
  \item[RQ3] How do environmental conditions affect library energy rankings,
    and can \gls{metrion-profiler} provide consistent measurements across
    them?
\end{description}

\section{Contributions}
\label{sec:intro-contributions}

This report extends the authors' 9th-semester
project~\cite{geotoms2025microbenchmarking}, which first integrated \gls{rapl}
measurement into BenchmarkDotNet and compared the libraries across a small set
of real-world payloads. The energy diagnoser is carried forward and reworked
here; the synthetic workload generator, the isolation and factorial designs,
and the \gls{metrion-profiler} validation are new to this report. With that
distinction in mind, the contributions are as follows.

\begin{itemize}
  \item A synthetic JSON workload generator that produces documents with
    controlled taxonomy dimensions --- size, nesting depth, structural width,
    content type, string composition, numeric representation, and redundancy.
  \item An energy-measurement extension to BenchmarkDotNet that exposes the
    Intel \gls{rapl} Package, Core, and DRAM domains as a diagnoser usable
    inside ordinary benchmarks.
  \item A two-part benchmark methodology --- single-dimension isolation
    sweeps and a byte-normalised factorial design --- that separates
    per-dimension energy scaling from confounding workload size.
  \item An empirical characterisation of the energy rankings and scaling
    rates of the five libraries across the workload dimensions (RQ1).
  \item An architectural and runtime account of the observed differences,
    drawing on allocation volume, garbage-collection behaviour, and
    source-level inspection (RQ2).
  \item A validation of the \gls{metrion-profiler} against \gls{rapl} and an
    assessment of ranking stability under controlled environmental stress
    (RQ3).
\end{itemize}

\section{Report Structure}
\label{sec:intro-structure}

The remainder of this report is organised as follows.
Chapter~\ref{chap:RelatedWork} reviews related work on software energy
measurement, JSON workload characterisation, and benchmark environment
control.
Chapter~\ref{chap:methodology} describes the experimental methodology: the
workload taxonomy, the generator, and the measurement design.
Chapter~\ref{chap:Implementation} details the energy-instrumented benchmark
harness and the libraries under test.
Chapter~\ref{chap:Results} reports the empirical results for RQ1--RQ3.
Chapter~\ref{chap:Discussion} draws the findings together into practical
implications, reflects on the design decisions, states the threats to
validity, and outlines future work.
Chapter~\ref{chap:Conclusion} concludes.


% =====================================================================
% NOTE: the Conclusion belongs at the END of the report, after the
% Discussion chapter. Kept here only for drafting.
%
% This is a FULLER recap (per request). The per-RQ answers are developed in
% \S\ref{sec:discussion-implications}; this chapter restates them in summary
% form. Watch for overlap with §6.1 when adjusting --- trim whichever side
% you want to carry the weight.
% =====================================================================

\chapter{Conclusion}
\label{chap:Conclusion}

This report investigated the energy consumption of five .NET JSON libraries
--- \texttt{System.Text.Json} in its reflection-based and source-generated
forms, \texttt{Newtonsoft.Json}, \texttt{SpanJson}, and \texttt{Utf8Json} ---
with three aims: to measure how their energy use depends on the shape of the
workload, to explain the differences between them, and to establish whether
those measurements can be trusted across environments. The study combined a
synthetic workload generator, an energy-instrumented extension of
BenchmarkDotNet exposing the Intel \gls{rapl} Package, Core, and DRAM
domains, and a two-part benchmark design: single-dimension isolation sweeps
that vary one workload property at a time, and a byte-normalised factorial
matrix that varies depth, width, and content type together. All benchmarks
were run on both a clean machine and a controlled-stress environment.

For RQ1, the energy rankings proved largely stable but not uniform.
\texttt{SpanJson} recorded the lowest deserialisation energy across most of
the workload space, and \texttt{STJSrcGen} the lowest serialisation energy in
the per-dimension analysis, while \texttt{Newtonsoft.Json} recorded the
highest energy in both operations almost everywhere. The orderings did shift
within individual sweeps. During deserialisation, string values longer than
fifty characters moved the lead to the \texttt{System.Text.Json} variants ---
\texttt{SpanJson}'s ratio to the lowest-energy library rose from $1.01\times$
at fifty characters to over $5\times$ at ten thousand --- and very wide
objects moved it to \texttt{Utf8Json}. The gap between the lowest- and
highest-energy libraries ranged from roughly $2.4\times$ on textual
deserialisation to over $11\times$ on Boolean serialisation, and the scaling
analysis attributed these shifts to differing per-unit growth rates rather
than to measurement noise.

For RQ2, the energy ordering tracked the libraries' execution-time ordering,
so the differences were explained through the CPU work that drives both ---
instruction count, allocation volume, and memory traffic --- rather than
through power draw, which the single-machine setup does not isolate. The
serialisation gap between the two \texttt{System.Text.Json} variants follows
from source generation, which removes at compile time the reflection that the
reflection-based variant performs at run time. Differences between
\texttt{SpanJson} and \texttt{System.Text.Json} were traced at the source
level to how each reads and writes values, with allocation volume held as a
control. \texttt{Newtonsoft.Json}'s consistently high energy, by contrast,
was not traced to a single mechanism and remains an open question.

For RQ3, the \gls{metrion-profiler}'s energy rankings agreed with those
derived from the hardware \gls{rapl} interface, and they remained stable when
the environment was placed under controlled stress, indicating that
consistent comparative measurements are achievable outside a quiet laboratory
setting. % TODO insert the concrete agreement measure (e.g. Spearman rank correlation) and the stress conditions tested
As a profiler still under development, its absolute readings and measurement
overhead carry the caveats set out in \S\ref{sec:threats}.

Taken together, these results support a practical guideline for energy-aware
library selection. Without workload-specific knowledge, \texttt{STJSrcGen} is
a sound default where serialisation dominates and \texttt{SpanJson} where
deserialisation of short string and primitive values dominates, while
\texttt{Newtonsoft.Json} is the highest-energy choice across the board, so a
migration away from it offers the largest single reduction. Where the
workload is known in advance, a more targeted choice is possible: applications
handling long textual payloads should prefer the \texttt{System.Text.Json}
variants, and those serialising very wide objects should consider
\texttt{Utf8Json}. Because several of the differences narrow to near-ties on
particular workloads, the size of the available reduction, rather than the
rank alone, should decide whether a change is worth making.

Beyond these findings, the study produced reusable artefacts --- the
synthetic workload generator, the \gls{rapl} energy diagnoser for
BenchmarkDotNet, and the validation harness for the \gls{metrion-profiler}
--- each of which can be applied to further energy studies, of JSON libraries
or otherwise.

The open questions raised here --- a fuller architectural account of
\texttt{Newtonsoft.Json}'s energy use, more realistic and byte-matched
workloads, and broader environmental and platform coverage --- are set out in
\S\ref{sec:future-work}. Measured directly and under controlled conditions,
energy is a tractable criterion for choosing between functionally equivalent
libraries, and this report shows that the choice can be made on evidence
rather than assumption.




SHort: 

      150 -This report set out to characterise the energy consumption of five .NET                                                                                                                          
      151 -JSON libraries, to explain the differences between them, and to validate what                                                                                                                    
      152 -it measured. To do so it combined a synthetic workload generator, an                                                                                                                             
      153 -energy-instrumented extension of BenchmarkDotNet exposing the Intel                                                                                                                              
      154 -\gls{rapl} domains, and a two-part benchmark design of single-dimension                                                                                                                          
      155 -isolation sweeps and a byte-normalised factorial matrix, run across both a                                                                                                                       
      156 -clean and a controlled-stress environment. 

      158 -The detailed answers to the three research questions are drawn together in                                                                                                                       
      159 -\S\ref{sec:discussion-implications}. The practical headline is that the                                                                                                                          
      160 -rankings are stable enough to serve as defaults: across most of the workload                                                                                                                     
      161 -space, \texttt{STJSrcGen} records the lowest serialisation energy and                                                                                                                            
      162 -\texttt{SpanJson} the lowest deserialisation energy, while                                                                                                                                       
      163 -\texttt{Newtonsoft.Json} consistently records the highest energy in both                                                                                                                         
      164 -operations. These defaults are not uniform --- long string values move the                                                                                                                       
      165 -deserialisation lead to the \texttt{System.Text.Json} variants, very wide                                                                                                                        
      166 -objects move it to \texttt{Utf8Json}, and the right choice depends on whether                                                                                                                    
      167 -an application's hot path is dominated by serialisation or deserialisation.                                                                                                                      
      168 -The largest energy reductions therefore come from selecting a library against                                                                                                                    
      169 -a known workload rather than in the abstract. 

      171 -Beyond the rankings, the generator, the energy diagnoser, and the                                                                                                                                
      172 -\gls{metrion-profiler} validation are reusable artefacts for further energy                                                                                                                      
      173 -studies. The open questions they raise --- deeper architectural explanations,                                                                                                                    
      174 -more realistic and byte-matched workloads, and broader environmental and                                                                                                                         
      175 -platform coverage --- are set out in \S\ref{sec:future-work}. Measured                                                                                                                           
      176 -directly and under controlled conditions, energy is a tractable criterion for                                                                                                                    
      177 -choosing between functionally equivalent libraries, and this report shows that                                                                                                                   
      178 -such a choice can be made on evidence rather than assumption.   
